% Source-derived chapter 07: Exotic classes
% Original source: standard-conjectures-abelian-fourfolds
\chapter{Exotic classes}
\label{standard-conjectures-abelian-fourfolds-chapter-07}
\source{standard-conjectures-abelian-fourfolds}

\begin{remark}[Source section]
\label{standard-conjectures-abelian-fourfolds-chapter-07-source}
\source{standard-conjectures-abelian-fourfolds}
\source{standard-conjectures-abelian-fourfolds}
This chapter reproduces the corresponding section of the retrieved
LaTeX source, including its definitions, statements, examples, and
proofs. It has not yet been translated into Lean.
\notready
\end{remark}

% The source section title was: Exotic classes
\label{sectionexotic}
\source{standard-conjectures-abelian-fourfolds}
In this section we fix an abelian variety $A$  of dimension four over a finite field $k$ (and we fix an algebraic closure $\bar{k}$ of $k$). After choosing a CM-structure for $A$ (Definition \ref{definitionCM}), Proposition \ref{primadecomposizione}(\ref{decomposizioneQ}) constructs motives ${\mathcal{M}_I \in \CHM (k)_{\Q}}$ which are direct factors of the motive  $\mathfrak{h}^{4}(A)$. Some of them, that we will call \textit{exotic}, are essential in the proof of the standard conjecture of Hodge type for $A$. The main result of the section (Proposition \ref{exoticdim2}) tells us that they are of rank two (in the sense of Definition \ref{defrank}).
\begin{definition}
\source{standard-conjectures-abelian-fourfolds}
\notready\label{definizione}
\source{standard-conjectures-abelian-fourfolds}
Let  $\mathcal{M}_I \in \CHM (k)_{\Q}$ be a  direct factor of $\mathfrak{h}^{4}(A)$ as constructed in Proposition \ref{primadecomposizione}(\ref{decomposizioneQ}). The motive $\mathcal{M}_I $  is called exotic if the  space of algebraic classes $\Hom_{\NUM(k)_\Q}(\one, \mathcal{M}_I (2))$ is nonzero and it does not contain any  nonzero Lefschetz class (Definition \ref{defexotic}).

An element in the $\mathbb{Q}$-vector space $\Hom_{\NUM(k)_\Q}(\one, \mathcal{M}_I (2))$ will be called an exotic class.
\end{definition}
\begin{remark}
\source{standard-conjectures-abelian-fourfolds}
\notready\label{remarkuguale}
\source{standard-conjectures-abelian-fourfolds}
Notice that, by Proposition \ref{talvez}(\ref{generatoalgebrico}), any exotic motive verifies the equality 
\begin{equation} \label{eq:uguale}
\source{standard-conjectures-abelian-fourfolds}
\dim_\Q \Hom_{\NUM(k)_\Q}(\one, \mathcal{M}_I (2)) = \dim \mathcal{M}_I.
 \tag{$\#$}
\end{equation}
Examples of exotic classes on abelian fourfolds (and especially of those that cannot be lifted to characteristic zero) will be discussed in  Appendix \ref{geometricexamples}.
\end{remark}
\begin{proposition}
\source{standard-conjectures-abelian-fourfolds}
\notready\label{exoticdim2}
\source{standard-conjectures-abelian-fourfolds}
Suppose that $A$ has dimension four. Then there exist a finite extension $k'$ of $k$ and a  CM-structure for $A\times_k k'$ such that any exotic motive $\mathcal{M}_I \in\CHM (k')_{\Q}$  of $\mathfrak{h}^{4}(A\times_k k')$ has dimension two. 
%More precisely, we have the following.
%\begin{enumerate}
%\item\label{puntouno} After possibly a finite extension of $k$, the space of exotic classes has dimension zero, two or four. 
%\item\label{puntodue} If the dimension of the space of exotic classes is four (even after  extension to the algebraic closure of $k$), then $A$ is isogenous to the product of a supersingular elliptic curve and an abelian threefold (possibly after a finite extension of $k$). In this case, one can choose a CM-structure so that the space of exotic classes will be decomposed in two different exotic motives.
%\item\label{puntotre} Any motive $\mathcal{M}_I $ of dimension one cannot be exotic.
%\end{enumerate}
\end{proposition}

The proof is decomposed in a series of lemmas and will take the rest of the section. We first fix notation and assumptions that will be used for the lemmas below.
\begin{notation}
\source{standard-conjectures-abelian-fourfolds}
\notready\label{notnotnot}
\source{standard-conjectures-abelian-fourfolds}
Let $\Sigma$ be as in Notation \ref{NotationCM}. Consider the decomposition $\mathfrak{h}^1(A) =\bigoplus_{\sigma \in \Sigma} M_\sigma$ from Proposition \ref{decomposition}. (Recall that  Frobenius acts on each of the eight $ M_\sigma$, see Remark \ref{remfrob}.) Let us denote  $\alpha_\sigma \in \overline{\Q}$ the  eigenvalues for the action of  Frobenius on  $ M_\sigma$. We will denote by $\bar{\cdot}$ the action of  complex conjugation on $ \Sigma$ (or on the set of parts of $ \Sigma$, or on $\overline{\Q}$). 
 
\end{notation}
\begin{assumption}
\source{standard-conjectures-abelian-fourfolds}
\notready\label{assumassum}
\source{standard-conjectures-abelian-fourfolds}
We will suppose that all the algebraic classes $\mathcal{Z}^*_\num (A_{\bar{k}})_\Q$ are defined over $k$. \end{assumption}
\begin{remark}
\source{standard-conjectures-abelian-fourfolds}
\notready\label{assumptionlef}
\source{standard-conjectures-abelian-fourfolds}
Note that the assumption above always holds after a finite extension of $k$. Notice also that, under this assumption, a class which becomes Lefschetz after a finite extension of the base field must be already Leftschetz. Hence, a motive which is exotic over the base field will still be exotic after a finite extension.
\end{remark}

\begin{lemma}
\source{standard-conjectures-abelian-fourfolds}
\notready\label{musica}
\source{standard-conjectures-abelian-fourfolds}
Let $q$ be the cardinality of $k$ and let $\mathcal{M}_I$ be an exotic motive. Then we have  the relation 
\begin{equation} \label{eq:1}
\source{standard-conjectures-abelian-fourfolds}
\prod_{\sigma \in I}\alpha_\sigma =q^2.
\end{equation}
Moreover, we have the property 
\begin{equation} \label{eq:2}
\source{standard-conjectures-abelian-fourfolds}
 \alpha_\sigma \cdot \alpha_\tau \neq q, \, \, \forall \sigma \neq \tau, \,\, \sigma,\tau \in I.
\end{equation} 
\end{lemma}

\begin{proof}
By Proposition \ref{talvez}(\ref{generatoalgebrico}) the $\ell$-adic realization of $\mathcal{M}_I$ is spanned by algebraic cycles hence  Frobenius acts on it by multiplication by $q^2$.  On the other hand, by construction (see  Proposition \ref{primadecomposizione}(\ref{decomposizioneL})),  Frobenius acts on the line $ M_I\subset \mathcal{M}_I$ via the multiplication by $\prod_{\sigma \in I}\alpha_\sigma $. This gives (\ref{eq:1}).   

Suppose now that  (\ref{eq:2}) is not satisfied.    Then $(\ref{eq:1})$ would force the set    $\{\alpha_\sigma\}_{\sigma \in I}$ to be of the form  $\alpha,q/\alpha, \beta, q/\beta$.   Each of the pairs $(\alpha,q/\alpha)$ and $(\beta, q/\beta)$ would correspond to a Frobenius invariant class in $H^2_\ell (A) (1)$. As the Tate conjecture for divisors on abelian varieties is known \cite[Theorem 4]{Tate}, each of these pairs corresponds to the class of a divisor, hence   $M_I$ would contain a Lefschetz class.
\end{proof}

\begin{definition}
\source{standard-conjectures-abelian-fourfolds}
\notready
A subset $I$ of $\Sigma$ of cardinality $4$ verifying (\ref{eq:1}) and (\ref{eq:2}) is called an exotic subset.
\end{definition}

\begin{remark}
\source{standard-conjectures-abelian-fourfolds}
\notready\label{completare}
\source{standard-conjectures-abelian-fourfolds}
Lemma \ref{musica} implies that the dimension of the space of exotic classes is at most the number of exotic subsets. The  Tate conjecture for $A$ predicts that this inequality should actually be an equality.
\end{remark}

\begin{lemma}
\source{standard-conjectures-abelian-fourfolds}
\notready\label{lemmaeven} The following holds.
\source{standard-conjectures-abelian-fourfolds}
\begin{enumerate}
\item\label{partizione} Let $I$ be an exotic subset, then the pair $(I,\bar{I})$ forms a partition of $\Sigma$.
\source{standard-conjectures-abelian-fourfolds}
\item Any  $\Gal(\overline{\Q}/\Q)$-conjugate of an exotic subset is again exotic.
\item Any exotic motive is even dimensional.
  \end{enumerate}
\end{lemma}
\begin{proof}
 Recall that the Weil conjectures imply $\bar{\alpha}=q/\alpha$. 
Then property (\ref{eq:2}) gives the first part. The second part follows from the very definition of exotic subset. 
Altogether complex conjugation acts without fixed points on  the $\Gal(\overline{\Q}/\Q)$-conjugates of $I$, hence the cardinality of the orbit of $I$ under the action of   $\Gal(\overline{\Q}/\Q)$ is even. On the other hand, this cardinality is the dimension of $\mathcal{M}_I$  (see  Proposition \ref{primadecomposizione}), this concludes the last part.
\end{proof}

\begin{lemma}
\source{standard-conjectures-abelian-fourfolds}
\notready
Under the Assumption \ref{assumassum}, an exotic subset $I$ verifies
\begin{equation} \label{eq:2n}
\source{standard-conjectures-abelian-fourfolds}
(\alpha_\sigma \cdot \alpha_\tau)^n \neq q^n, \, \, \forall \sigma \neq \tau, \,\, \sigma,\tau \in I
\end{equation} 
(for all positive integer $n$).
\end{lemma}
\begin{proof}
 If we extend the scalars to  $\mathbb{F}_{q^n}$  the motive $M_I$  will still be exotic (Remark \ref{assumptionlef}).  The eigenvalues of  Frobenius become $\{\alpha_\sigma^n\}$. Hence, by applying  (\ref{eq:2}) over this new base field we deduce (\ref{eq:2n}).
\end{proof}
\begin{lemma}
\source{standard-conjectures-abelian-fourfolds}
\notready\label{lemmanotwo}
\source{standard-conjectures-abelian-fourfolds}
Under the Assumption \ref{assumassum}, two exotic subsets cannot intersect in exactly two elements.
\end{lemma}
\begin{proof}
 Suppose that there are two exotic subsets  $I$ and $J$ whose intersection has cardinality two. Let us call $\alpha,\beta$ the two Frobenius eigenvalues corresponding to  $I-I\cap J$ and $\gamma,\delta$  those corresponding to  $J-I\cap J$.  Property (\ref{eq:1})  gives
 \[\alpha\cdot \beta=\gamma \cdot  \delta\]
and  property (\ref{eq:2}) gives (after reordering if necessary)
\[\alpha= q/ \gamma \,\, , \,\, \beta=  q/\delta.\]
Putting these relations together we have 
$\alpha^2 =q^2/\beta^2.$ This equality gives a contradiction by applying   (\ref{eq:2n})   to $I$ for $n=2$.
\end{proof}



\begin{lemma}
\source{standard-conjectures-abelian-fourfolds}
\notready\label{lemmaabove}
\source{standard-conjectures-abelian-fourfolds}
Under the Assumption \ref{assumassum}, the dimension of the space of exotic classes is zero, two or four. More precisely, there are either  zero, two or four exotic subsets. In case they are two, they are  complex conjugate to each other. In case they are four, they are of the form $I,\bar{I}, J, \bar{J}$, with $I\cap J$ of cardinality three.
 \end{lemma}
\begin{proof}
The statement on exotic classes is implied by the statement on exotic subsets. Indeed, the dimension of the space of exotic classes is even (Lemma \ref{lemmaeven} together with the relation (\ref{eq:uguale}) in Remark \ref{remarkuguale}) and it is  at most the number of exotic subsets (Remark \ref{completare}).


Let us now show the statement on exotic subsets. We have already pointed out in the proof of Lemma \ref{lemmaeven} that  complex conjugation acts without fixed points on the exotic subsets. Suppose now that there are at least four such subsets, call them $I,\bar{I}, J, \bar{J}$. Then one subset among $J$ and $\bar{J}$ intersects $I$ in at least two elements. Without loss of generality we suppose it is $J$. Then by Lemma \ref{lemmanotwo}, $I\cap J$ must be of cardinality three.

If there were more than four exotic subsets, then, by the same arguments there would be an exotic subset $K$ intersecting $I$ in exactly three elements. Then the intersection of $J$ and $K$ would have exactly two elements. This is impossible  by Lemma \ref{lemmanotwo}.
\end{proof}

\begin{lemma}
\source{standard-conjectures-abelian-fourfolds}
\notready\label{lemmafour}
\source{standard-conjectures-abelian-fourfolds}
Suppose that there are four exotic subsets and that the Assumption   \ref{assumassum} is satisfied. Then, after extending $k= \mathbb{F}_{q}$ to its quadratic extension $\mathbb{F}_{q^2}$, the abelian fourfold $A$ becomes isogenous to  $E\times X$ where  $X$ is an abelian threefold and $E$ is a supersingular elliptic curve on which   Frobenius acts as $q\cdot \id $. 


\end{lemma}
 
\begin{proof}
We keep notation from Lemma \ref{lemmaabove}. If the element of $I-I\cap J$ is $\sigma$ then the element of  $(J-I\cap J)= J \cap \bar{I}$ must be $\bar{\sigma}$ because of  Lemma \ref{lemmaeven}(\ref{partizione}).  Let $\alpha$ be the Frobenius eigenvalue for the action on $M_\sigma$, then  (\ref{eq:1}) applied to $I$ and $J$ implies $\alpha=\bar{\alpha}$, hence $\alpha^2=q$.

Let us now extend the field of definition to $\mathbb{F}_{q^2}$. Among the eight eigenvalues of  Frobenius we will find $q$, which means that there is a nonzero Frobenius-equivariant map between the Tate module of the supersingular elliptic curve $E$ and the Tate module of $A$. This implies the statement by  \cite[Theorem 4]{Tate}.
\end{proof}


\begin{lemma}
\source{standard-conjectures-abelian-fourfolds}
\notready
We keep Notation \ref{notnotnot} and Assumption \ref{assumassum}. Consider on the base field $k=\mathbb{F}_{q^2}$   a supersingular elliptic curve $E$  on which   Frobenius acts as $q\cdot \id$.
Let $A$ be an abelian fourfold of the form $ A=E\times X$, let  $\Sigma_X$ and $ \Sigma_E$ be CM-structures for $X$ and $E$ and define  $\Sigma=\Sigma_X \cup \Sigma_E$. Then, any exotic subset $I$ verifies that $I\cap \Sigma_X$ has cardinality three. Moreover,  the motive $\mathcal{M}_{I\cap \Sigma_X} \in\CHM (k)_{\Q}$, direct factor of $\mathfrak{h}^3(X)$ as constructed in Proposition \ref{primadecomposizione}, is of dimension two. 

Finally, if the space of exotic classes on $A$ is four dimensional, it is contained in the four dimensional motive $\mathcal{M}_{I\cap \Sigma_X} \otimes \mathfrak{h}^1(E)$.
\end{lemma}
\begin{proof}
First recall that  the pair $(I,\bar{I})$ forms a partition of $\Sigma$ by Lemma \ref{lemmaeven}(\ref{partizione}). As $ \Sigma_X$ has cardinality six and is stable by the action of complex conjugation one must have that $I\cap \Sigma_X$ has cardinality three.


Consider the subsets $K \subset \Sigma_X$ (of cardinality three) verifying
\begin{equation} \label{eq:4}
\source{standard-conjectures-abelian-fourfolds}
\prod_{\sigma \in K} \alpha_\sigma =q^3
\end{equation}
\begin{equation} \label{eq:5}
\source{standard-conjectures-abelian-fourfolds}
 \alpha_\sigma \cdot \alpha_\tau \neq q^2, \, \, \forall \sigma \neq \tau, \,\, \sigma,\tau \in K
\end{equation}
and
\begin{equation} \label{eq:5epoi}
\source{standard-conjectures-abelian-fourfolds}
 \alpha_\sigma  \neq q,  \,\, \forall \sigma\in K
\end{equation}

Clearly the relations (\ref{eq:1}) and (\ref{eq:2}) for $I$ imply that $K=I\cap \Sigma_B$ verifies (\ref{eq:4}),(\ref{eq:5}) and (\ref{eq:5epoi}). Conversely, the relations (\ref{eq:4}),(\ref{eq:5}) and (\ref{eq:5epoi})  for $K$ imply that $I=K\cup \{\sigma\}$ verifies (\ref{eq:1}) and (\ref{eq:2}) for any $\sigma \in \Sigma_E$. As there are at most four exotic subsets for $A$, there must be at most two subsets of $\Sigma_X$ verifying  (\ref{eq:4}),(\ref{eq:5}) and (\ref{eq:5epoi}).  On the other hand $I\cap \Sigma_X$  and $\bar{I}\cap \Sigma_X$ are two of those, which implies that there are exactly  two subsets of $\Sigma_X$  verifying (\ref{eq:4}),(\ref{eq:5}) and (\ref{eq:5epoi}).
Finally, as the Galois group $\Gal(\bar{\Q}/\Q)$ acts on those subsets then the motive $\mathcal{M}_{I\cap \Sigma_X}$ has dimension two. 
The rest follows from the construction of $\mathcal{M}_{I\cap \Sigma_X}$.
\end{proof}
\begin{lemma}
\source{standard-conjectures-abelian-fourfolds}
\notready\label{bellaidea}
\source{standard-conjectures-abelian-fourfolds}
Consider an abelian fourfold of the form $ A=E\times X$, where  $E$ is a supersingular elliptic curve such that $\dim_\Q\End(E)\otimes_\Z \Q=4$. Fix a CM-structure for $X$ and consider a direct factor of $\mathfrak{h}^3(X)$  of the form $\mathcal{M}_{I} \in\CHM (k)_{\Q}$, as constructed in Proposition \ref{primadecomposizione}. Suppose that $\mathcal{M}_{I} $ has rank two and that  the  motive ${\mathcal{M}_I \otimes \mathfrak{h}^1(E)}$ is spanned by algebraic classes. Then  there exists a CM structure for  $E$ (write $\Sigma_E=\{\sigma,\bar{\sigma}\}$) such that  the motive $\mathcal{M}_I \otimes \mathfrak{h}^1(E)$ decomposes in the category $\CHM (k)_{\Q}$ as the sum of two motives of rank two 
\[\mathcal{M}_I \otimes \mathfrak{h}^1(E)= \mathcal{M}_{I\cup \sigma} \oplus \mathcal{M}_{I\cup \bar{\sigma}}\]
via Proposition  \ref{primadecomposizione}(\ref{decomposizioneQ}).
 \end{lemma}
\begin{proof}

By construction of   $\mathcal{M}_I$ (Proposition  \ref{primadecomposizione}(\ref{decomposizioneQ})) we can find a quadratic number field $F$  such that the motive $\mathcal{M}_{I} $  decomposes in the category $\CHM (k)_{F}$ into the sum 
\[\mathcal{M}_{I} = M_{I} \oplus  M_{\bar{I}}  \] of two motives of rank one. In order to conclude, it is enough to show that also $\mathfrak{h}^1(E)$  decomposes in the category $\CHM (k)_{F}$ into the sum of two motives of dimension one. 

This fact is equivalent to saying that $F$ can be embedded in the division algebra $\End(E)\otimes_\Z \Q.$ As this algebra splits at all primes except infinity and  the characteristic of the base field (call it $p$), we have to show that the embeddings $F \subset \R$ and $F\subset \Q_p$ cannot exist. Indeed, as the degree of $F$ over $\Q$ is the same as the index of the central $\Q$-algebra $\End(E)\otimes_\Z \Q$ (namely two), the inclusion  $F\subset\End(E)\otimes_\Z \Q$ is equivalent to the fact that $\End(E)\otimes_\Z F$ is a matrix algebra. By the Albert--Brauer--Hasse--Noether theorem this can be checked locally. As $\End(E)\otimes_\Z \Q_\ell$ is already a matrix algebra it is enough to check at $p$ and at infinity. Now if $F $ does not embed in $\R$ and $\Q_p$ then $F\otimes_{\Q} \R$ and $F\otimes_{\Q}\Q_p$ are quadratic field extensions of $\R$ and $\Q_p$. As $\End(E)\otimes_\Z \Q$ has invariant $1/2$ at $p$ and infinity it becomes a matrix algebra after tensoring with these quadratic extensions.

We will give two proofs of the non existence of the embeddings $F \subset \R$ and $F\subset \Q_p$.

A first proof uses that these motives lift to characteristic zero (Proposition \ref{lifting motives}). 
Now if $F$ was contained in $\R$ then the Betti realization of the lifting of $M_{I}$ would respect the Hodge symmetry. This is impossible as it is one dimensional and of weight three. Similarly, if $F$ was contained in $\Q_p$ then the Hyodo-Kato realization of $M_{I}$ and $M_{\bar{I}} $ would be two filtered $\varphi$-modules of dimension one with different filtrations (again because the weight of $\mathcal{M}_{I}$ is three). This implies that  
$\det (M_{I}  \otimes \mathfrak{h}^1(E))$ and $\det(M_{\bar{I}}  \otimes \mathfrak{h}^1(E))$ would realize in two filtered $\varphi$-modules of dimension one and different filtrations.  On the other hand   absolute Frobenius acts on both   in the same way (namely by multiplication by $p^4$) because they are spanned by algebraic classes. This implies that at least one of the two filtered $\varphi$-modules is not admissible and concludes the proof as all filtered $\varphi$-modules coming from geometry (in particular filtered $\varphi$-modules that are realizations of motives) must be admissible.

We give an alternative proof, which does not use the lifting to characteristic zero. We write it for $\Q_p$, it works in the same way for $\R$. If $F$ was contained in $\Q_p$ then the motive $M_{I}  \otimes \mathfrak{h}^1(E)$  would live in $\CHM (k)_{\Q_p}$. On such a motive the division algebra $\End(E)\otimes_\Z \Q_p$ acts hence it acts also on the $\Q_p$ vector space spanned by the algebraic classes of the motive. On the other hand, as this motive is spanned by algebraic classes, the division algebra  $\End(E)\otimes_\Z \Q_p$  would act on a $\Q_p$-vector space  of dimension two. This gives a contradiction as such an  action cannot exist.
\end{proof}
\begin{proof}[Proof of Proposition \ref{exoticdim2}] This is a combination of the  last four lemmas.
\end{proof}
